\chapter{Summary and Outlook}
\label{c:summary}

%%%%%%%%%%%%%%%%%%%%%%%%%%%%%%%%%%%%%%%%%%%%%%%%%%%%%%%%%%%%%%%%%%%%%%%%%%%%%%%%

% Summary

This thesis presented the development of an efficient, matrix-free C++ simulator for second-quantized fermionic operators.
To overcome the exponential memory scaling associated with tracking identical particles in many-body quantum systems, the simulator implements a sparse state representation.
By utilizing an unordered hash map the occupation number representation of the Fock space, can be implemented as linking unsigned integer bitstrings to complex probability amplitudes, which avoids saving zero-amplitude states.
Furthermore, the system encodes arbitrary fermionic Generators and Hamiltonian's by storing their associated creation and annihilation operators using integer vectors.
This ensures that all possible combinations of operators can be stored and therefore calculated.
Also, we can therefore implement a straightforward string representations like  Google's OpenFermion.

A central achievement of this work is the implementation of matrix-free direct state evolution.
Because constructing global unitary matrices for large Hilbert spaces is memory prohibitive, the simulator applies excitation generators directly to the computational basis.
The algorithm evaluates fermionic parity phases and anticommutation relations by tracking orbital occupancies and utilizing native CPU population counts.
This eliminates the need for Jordan-Wigner strings and enforces the strict antisymmetry required by the Pauli exclusion principle.

Beyond basic excitations, the simulator calculates fermionic SWAP operations and expectation value evaluations by iterating over non-zero amplitudes.
It uses bitwise manipulations, such as XOR operations and bit-shifting for nullspace projection, in order to save memory and speed up the operation.


% Outlook
Although the core algorithmic design and C++ implementation provides a foundation, several extensions can potentially further increase its performance and practical utility.
A natural step for future development could focus on parallelization and integration with quantum chemistry ecosystems.

First, the simulator’s core operations are suited for parallel execution.
Currently, operations such as applying fermionic excitations and evaluating expectation values iterate over the basis states within the hash map.
Because each basis state update and nullspace check can be evaluated independently, these inner loops can be parallelized across multiple CPU threads.
Similar with calculating expectation values.
Here large sums of Hamiltonian terms are getting calculated one by one which could be parallelized as well, which might decrease runtimes for complex molecular geometries.
Another not as obvious speedup, could be to vectorize the bitwise state manipulations, mask evaluations, and population counts (e.g. using GPU operations) to process large batches of simultaneously.

Beyond optimization, practical application requires integration into higher level workflows.
To bridge the gap between C++ performance and ease of use, an ongoing extension is the development of a dedicated Python wrapper module.
By exposing the underlying operations to Python, the simulator can be integrated into quantum chemistry packages.

The following features could be implemented with such a wrapper.
It enables the simulator to serve as an expectation-value backend for parameterized fermionic circuits.
Users can supply molecular properties, such as one- and two-body integrals, orbital occupations, and electronic spin and build the full fermionic Hamiltonian in Python.
The wrapper also manages orbital reordering between different spin representations and internal up-then-down formats, maps string operators into native FermionTerm instances, and evaluates parameterized states without exposing the underlying engine.
Finally, this enables the simulator to be used within Variational Quantum Eigensolver (VQE) algorithms.
This combines the execution speed of C++ operations with the flexibility of computing platforms.

